\section{Entity-Aware Evaluation}
\label{sec:evaluation}

\subsection{Evaluation Framework}
\label{sec:evaluation-framework}

Evaluation follows the structure and semantics of the target schema. Different
metadata fields require different evaluation metrics: constrained values use
exact or normalized-exact match, free text and wording-sensitive names use
ROUGE-L, and unordered lists use set F1. The metric columns in
Tables~\ref{tab:appendix-iaa} and~\ref{tab:appendix-zero-shot} identify the
metric applied to every scored schema field.

Paper-level fields can be evaluated directly against their gold values,
whereas entity collections first require correspondence between
predicted and gold entities. Our evaluation therefore separates three
outcomes: entity recovery, attribute quality among aligned entities, and
reference quality among aligned entities. Because these outcomes
have different meanings and supports, we report them separately rather than
constructing a schema-wide average.


\subsection{Entity Alignment}
\label{sec:entity-alignment}

The six entity collections require one-to-one alignment:
affiliations, authors, funding, facilities, samples, and datasets.

For authors and affiliations, each record provides a direct textual alignment
key: \texttt{full\_name} for authors and \texttt{text} for affiliations.
After normalization, predicted--gold pairs with token-overlap F1 of at least
0.5 are considered candidate matches. Hungarian assignment then selects a
one-to-one matching that maximizes the number of correspondences, with total
similarity used to break ties.

Facilities, funding, samples, and datasets require collection-specific
alignment because their identifying fields can vary in wording and because
similar records may correspond to distinct entities rather than the same one.
For these collections, the evaluator uses the following fields as the primary
basis for alignment: facility \texttt{name}; funding \texttt{grant\_id} and
\texttt{funder}; sample \texttt{sample\_name} and
\texttt{chemical\_composition}; and dataset \texttt{dataset\_name} and
\texttt{measurement\_technique}. Where necessary, these fields are
supplemented with reviewed alternative forms grounded in the source PDF and
entity distinctions recorded during reconciliation. Alternative forms allow
different wording to denote the same entity, whereas entity distinctions
prevent related but separate records from being collapsed. The original field
values are retained for subsequent attribute scoring.

The matching criteria are applied in collection-specific stages. Facilities
are matched by exact name, then by reviewed alternative names for the same
facility, and then by unambiguous complete name containment. For example,
``ESRF'' can correspond to ``European Synchrotron Radiation Facility'' when
the paper defines the acronym, whereas an instrument or beamline named without 
its host facility does not match the facility. Funding records are matched
primarily by normalized grant identifier, with funder information used to 
distinguish records where needed.

Samples are aligned by their names and compositions in the same staged way: first an exact
sample-name match; then a composition that matches and uniquely identifies a
gold sample within the paper; then, among samples of the same composition, 
a match based on distinguishing name words, such as ``with'' sonication versus 
``without'' sonication, or explicit values such as temperature, thickness, or concentration; 
finally, remaining samples are matched when one name contains all words
of the other, as in ``MnSi'' and ``MnSi single crystal'', unless their
reported compositions disagree or another gold sample is an equally
supported match. 

Datasets are aligned similarly using their names and measurement techniques: 
first by the same name with a compatible measurement technique; then by a name that uniquely 
identifies a gold dataset within the paper; then, among datasets sharing a technique, by 
distinguishing words in their names; and finally by technique alone, when the technique identifies 
a single gold dataset and matches only one remaining eligible prediction.

Within each collection, later stages are applied only to entities left
unmatched by earlier stages. Matches follow the collection-specific
criteria above and preserve one-to-one correspondence.

\subsection{Scoring and Aggregation}
\label{sec:scoring-aggregation}

Paper-level fields are scored directly against their gold values using the
field-specific metrics listed in Tables~\ref{tab:appendix-iaa}
and~\ref{tab:appendix-zero-shot}. Scalar paper-level scores are averaged
across papers, while count-based paper-level F1 metrics are micro-aggregated
by pooling true positives, false positives, and false negatives before
computing F1. Section-heading content is scored with set F1; the relative order
of headings is reported only as an appendix diagnostic using Kendall's
$\tau$.

For entity collections, the one-to-one alignment
(Section~\ref{sec:entity-alignment}) determines entity recovery. Aligned
predicted--gold pairs count as true positives, unmatched predictions as false
positives, and unmatched gold entities as false negatives. Entity F1 is
micro-aggregated across the corpus by pooling these counts before computing
F1, so entity-dense papers contribute proportionally.

Attribute metrics are applied only to aligned entity pairs and therefore
measure attribute quality conditional on entity recovery. Scalar attribute
scores are averaged over aligned pairs, while count-based attribute F1 metrics
are micro-aggregated by pooling counts across aligned pairs and papers. Because
the aligned set depends on recovery, conditional attribute scores should be
interpreted together with the corresponding collection's entity F1.

Cross-entity references are scored after predicted output-local identifiers
are remapped through the established entity alignments. For example, if
predicted \texttt{sample\_7} aligns with gold \texttt{sample\_2}, a reference
to \texttt{sample\_7} is evaluated as a reference to
\texttt{sample\_2}. Reference F1 is micro-aggregated across the corpus by
pooling reference true positives, false positives, and false negatives over
aligned parent entities. Reference F1 is therefore conditional and should be
interpreted together with the corresponding parent collection's entity F1.
